\documentclass[11pt,a4paper,oneside]{book}

% ── Encoding & Fonts ──
\usepackage[T1]{fontenc}
\usepackage[utf8]{inputenc}
\usepackage{lmodern}
\usepackage{microtype}

% ── Page layout ──
\usepackage[margin=2.5cm,top=3cm,bottom=3cm]{geometry}

% ── Colors ──
\usepackage[dvipsnames,svgnames,x11names]{xcolor}

% Project color palette
\definecolor{pyblue}{HTML}{306998}
\definecolor{pyyellow}{HTML}{FFD43B}
\definecolor{goblue}{HTML}{00ADD8}
\definecolor{rustred}{HTML}{DEA584}
\definecolor{javaorange}{HTML}{F89820}
\definecolor{codebg}{HTML}{F5F5F0}
\definecolor{warnbg}{HTML}{FFF3CD}
\definecolor{warnborder}{HTML}{856404}
\definecolor{dangerbg}{HTML}{F8D7DA}
\definecolor{dangerborder}{HTML}{721C24}
\definecolor{infobg}{HTML}{D1ECF1}
\definecolor{infoborder}{HTML}{0C5460}
\definecolor{successbg}{HTML}{D4EDDA}
\definecolor{successborder}{HTML}{155724}
\definecolor{chapterbg}{HTML}{2B5B84}

% ── Code listings ──
\usepackage{listings}
\usepackage{inconsolata}

\lstdefinestyle{python}{
    language=Python,
    basicstyle=\ttfamily\small,
    backgroundcolor=\color{codebg},
    keywordstyle=\bfseries\color{pyblue},
    stringstyle=\color{OliveGreen},
    commentstyle=\itshape\color{Gray},
    numberstyle=\tiny\color{Gray},
    numbers=left,
    numbersep=8pt,
    frame=l,
    framerule=2pt,
    rulecolor=\color{pyblue},
    breaklines=true,
    breakatwhitespace=true,
    showstringspaces=false,
    tabsize=4,
    xleftmargin=16pt,
    framexleftmargin=16pt,
    aboveskip=12pt,
    belowskip=8pt,
    morekeywords={async,await,match,case,TaskGroup,ExceptionGroup},
}

\lstdefinestyle{go}{
    language=Go,
    basicstyle=\ttfamily\small,
    backgroundcolor=\color{codebg},
    keywordstyle=\bfseries\color{goblue},
    stringstyle=\color{OliveGreen},
    commentstyle=\itshape\color{Gray},
    frame=l,
    framerule=2pt,
    rulecolor=\color{goblue},
    breaklines=true,
    showstringspaces=false,
    tabsize=4,
    xleftmargin=16pt,
    framexleftmargin=16pt,
    aboveskip=12pt,
    belowskip=8pt,
}

\lstdefinestyle{rust}{
    language=C, % approximate
    basicstyle=\ttfamily\small,
    backgroundcolor=\color{codebg},
    keywordstyle=\bfseries\color{BrickRed},
    stringstyle=\color{OliveGreen},
    commentstyle=\itshape\color{Gray},
    frame=l,
    framerule=2pt,
    rulecolor=\color{rustred},
    breaklines=true,
    showstringspaces=false,
    tabsize=4,
    xleftmargin=16pt,
    framexleftmargin=16pt,
    aboveskip=12pt,
    belowskip=8pt,
    morekeywords={fn,let,mut,async,await,pub,impl,struct,enum,trait,use,mod,tokio,spawn},
}

\lstdefinestyle{java}{
    language=Java,
    basicstyle=\ttfamily\small,
    backgroundcolor=\color{codebg},
    keywordstyle=\bfseries\color{javaorange!80!black},
    stringstyle=\color{OliveGreen},
    commentstyle=\itshape\color{Gray},
    frame=l,
    framerule=2pt,
    rulecolor=\color{javaorange},
    breaklines=true,
    showstringspaces=false,
    tabsize=4,
    xleftmargin=16pt,
    framexleftmargin=16pt,
    aboveskip=12pt,
    belowskip=8pt,
    morekeywords={var,record,sealed,permits,yield},
}

\lstdefinestyle{toml}{
    basicstyle=\ttfamily\small,
    backgroundcolor=\color{codebg},
    keywordstyle=\bfseries\color{pyblue},
    stringstyle=\color{OliveGreen},
    commentstyle=\itshape\color{Gray},
    frame=l,
    framerule=2pt,
    rulecolor=\color{Gray},
    breaklines=true,
    showstringspaces=false,
    tabsize=4,
    xleftmargin=16pt,
    framexleftmargin=16pt,
    aboveskip=12pt,
    belowskip=8pt,
}

\lstdefinestyle{shell}{
    basicstyle=\ttfamily\small,
    backgroundcolor=\color{codebg},
    keywordstyle=\bfseries\color{pyblue},
    frame=l,
    framerule=2pt,
    rulecolor=\color{DarkGreen},
    breaklines=true,
    showstringspaces=false,
    xleftmargin=16pt,
    framexleftmargin=16pt,
    aboveskip=12pt,
    belowskip=8pt,
}

\lstdefinestyle{directory}{
    basicstyle=\ttfamily\small,
    backgroundcolor=\color{codebg},
    frame=single,
    rulecolor=\color{Gray},
    breaklines=true,
    showstringspaces=false,
    xleftmargin=16pt,
    framexleftmargin=16pt,
    aboveskip=12pt,
    belowskip=8pt,
}

\lstset{style=python}

% ── Boxes and admonitions ──
\usepackage[most]{tcolorbox}

\newtcolorbox{warningbox}[1][]{
    colback=warnbg,
    colframe=warnborder,
    coltitle=warnborder,
    fonttitle=\bfseries,
    title={Warning},
    arc=2pt,
    boxrule=1pt,
    left=8pt,
    right=8pt,
    #1,
}

\newtcolorbox{dangerbox}[1][]{
    colback=dangerbg,
    colframe=dangerborder,
    coltitle=dangerborder,
    fonttitle=\bfseries,
    title={Do NOT},
    arc=2pt,
    boxrule=1pt,
    left=8pt,
    right=8pt,
    #1,
}

\newtcolorbox{infobox}[1][]{
    colback=infobg,
    colframe=infoborder,
    coltitle=infoborder,
    fonttitle=\bfseries,
    title={Note},
    arc=2pt,
    boxrule=1pt,
    left=8pt,
    right=8pt,
    #1,
}

\newtcolorbox{successbox}[1][]{
    colback=successbg,
    colframe=successborder,
    coltitle=successborder,
    fonttitle=\bfseries,
    title={Best Practice},
    arc=2pt,
    boxrule=1pt,
    left=8pt,
    right=8pt,
    #1,
}

\newtcolorbox{analogybox}[1][]{
    colback=white,
    colframe=pyblue!50!black,
    coltitle=white,
    colbacktitle=pyblue!70!black,
    fonttitle=\bfseries\sffamily,
    title={Cross-Language Analogies},
    arc=3pt,
    boxrule=1pt,
    left=8pt,
    right=8pt,
    #1,
}

\newtcolorbox{migrationbox}[1][]{
    colback=pyyellow!10!white,
    colframe=pyyellow!60!black,
    coltitle=pyyellow!20!black,
    colbacktitle=pyyellow!80,
    fonttitle=\bfseries\sffamily,
    title={Migration from older Python},
    arc=3pt,
    boxrule=1pt,
    left=8pt,
    right=8pt,
    #1,
}

% ── Chapter/section styling ──
\usepackage{titlesec}
\usepackage{titletoc}

\titleformat{\chapter}[display]
    {\normalfont\huge\bfseries\sffamily\color{chapterbg}}
    {\chaptertitlename\ \thechapter}{20pt}{\Huge}
\titleformat{\section}
    {\normalfont\Large\bfseries\sffamily\color{pyblue!80!black}}
    {\thesection}{1em}{}
\titleformat{\subsection}
    {\normalfont\large\bfseries\sffamily\color{pyblue!60!black}}
    {\thesubsection}{1em}{}

% ── Hyperlinks ──
\usepackage[
    colorlinks=true,
    linkcolor=pyblue!80!black,
    urlcolor=pyblue,
    citecolor=pyblue,
    bookmarks=true,
    bookmarksnumbered=true,
    pdfauthor={Project Standards},
    pdftitle={Python 3.14 Async Backend: A Practical Guide},
    pdfsubject={asyncio, structured concurrency, Python 3.14},
]{hyperref}

% ── Tables ──
\usepackage{booktabs}
\usepackage{tabularx}
\usepackage{array}

% ── Graphics ──
\usepackage{graphicx}

% ── Header/footer ──
\usepackage{fancyhdr}
\pagestyle{fancy}
\fancyhf{}
\fancyhead[L]{\sffamily\small\leftmark}
\fancyhead[R]{\sffamily\small Python 3.14 Async Guide}
\fancyfoot[C]{\thepage}
\renewcommand{\headrulewidth}{0.4pt}
\renewcommand{\headrule}{\hbox to\headwidth{\color{pyblue}\leaders\hrule height \headrulewidth\hfill}}

% ── Misc ──
\usepackage{enumitem}
\setlist{nosep,leftmargin=*}

\newcommand{\py}[1]{\texttt{\color{pyblue}#1}}
\newcommand{\gocode}[1]{\texttt{\color{goblue}#1}}
\newcommand{\rustcode}[1]{\texttt{\color{BrickRed}#1}}
\newcommand{\javacode}[1]{\texttt{\color{javaorange!80!black}#1}}
\newcommand{\deprecated}[1]{\texttt{\color{red!70!black}\sout{#1}}}

\usepackage[normalem]{ulem} % for \sout

% ════════════════════════════════════════════════════════
\begin{document}

% ── Title page ──
\begin{titlepage}
\begin{center}
\vspace*{2cm}

{\fontsize{40}{48}\selectfont\bfseries\sffamily\color{chapterbg}
Python 3.14\\[4pt]
Async Backend}

\vspace{0.5cm}

{\fontsize{20}{24}\selectfont\sffamily\color{pyblue}
A Practical Guide}

\vspace{1.5cm}

\rule{0.6\textwidth}{1pt}

\vspace{1.5cm}

{\large\sffamily
Structured Concurrency, Modern Tooling,\\
and Production Patterns}

\vspace{2cm}

\begin{tcolorbox}[
    colback=codebg,
    colframe=pyblue,
    width=0.75\textwidth,
    arc=4pt,
    boxrule=1.5pt,
    halign=center,
]
{\sffamily\small
For senior engineers coming from\\[4pt]
\textcolor{goblue}{\textbf{Go}} $\cdot$
\textcolor{javaorange}{\textbf{Java/Kotlin}} $\cdot$
\textcolor{BrickRed}{\textbf{Rust}}\\[4pt]
who know async basics but are new to Python 3.14}
\end{tcolorbox}

\vfill

{\small\sffamily\color{Gray} March 2026}
\end{center}
\end{titlepage}

% ── TOC ──
\tableofcontents
\newpage

% ════════════════════════════════════════════════════════
\chapter{How Async I/O Works}

Before diving into Python-specific APIs, you need to understand \emph{what async actually does}
at the operating system level. If you've only used Go (where the runtime hides all of this),
this chapter fills that gap.

\section{The Problem: One Thread Per Connection}

Traditional blocking I/O assigns one OS thread per connection. When a thread calls
\texttt{read()} on a socket, the kernel \emph{suspends} that thread until data arrives.
The thread consumes memory (typically 1--8\,MB of stack) and does nothing while waiting.

At 10,000 concurrent connections, that's 10,000 threads --- tens of gigabytes of memory and
constant context-switching overhead. This is the \textbf{C10K problem} that async I/O solves.

\begin{analogybox}
\begin{itemize}
    \item \textcolor{goblue}{\textbf{Go:}} Goroutines solve the same problem but differently ---
          they use small stacks ($\sim$4\,KB, growable) and the runtime multiplexes them onto a small
          pool of OS threads via an M:N scheduler. You write blocking-style code, but it's
          non-blocking under the hood.
    \item \textcolor{javaorange}{\textbf{Java 21+:}} Virtual threads (Project Loom) take the Go
          approach --- lightweight threads scheduled by the JVM. Before Loom, Java used NIO with
          selectors (the same epoll/kqueue approach Python uses).
    \item \textcolor{BrickRed}{\textbf{Rust:}} Tokio is an explicit async runtime like Python's
          asyncio --- \texttt{async/await} with an event loop backed by epoll/kqueue/io\_uring.
\end{itemize}
\end{analogybox}

\section{Non-Blocking I/O and I/O Multiplexing}

The solution is two-part:

\begin{enumerate}
    \item \textbf{Non-blocking sockets.} Set \texttt{O\_NONBLOCK} on a file descriptor. Now
          \texttt{read()} returns immediately with \texttt{EAGAIN} if no data is available,
          instead of suspending the thread.
    \item \textbf{I/O multiplexing.} Ask the kernel: ``tell me when \emph{any} of these
          file descriptors have data ready.'' One thread monitors thousands of connections.
\end{enumerate}

The kernel provides several multiplexing APIs, each an evolution of the last:

\begin{center}
\renewcommand{\arraystretch}{1.3}
\begin{tcolorbox}[
    colback=codebg,
    colframe=pyblue,
    width=0.98\textwidth,
    arc=3pt,
    boxrule=1pt,
]
\small
\begin{tabularx}{\textwidth}{>{\bfseries\ttfamily}l l X}
\toprule
API & Platform & How it works \\
\midrule
select() & All & Pass a bitmap of fds. Kernel scans all of them. $O(n)$ per call.
           Limited to 1024 fds. Ancient, slow. \\
poll() & All & Like select but no fd limit. Still $O(n)$ --- kernel scans the entire list. \\
epoll & Linux & Register fds once. Kernel maintains an interest list internally and
         returns \emph{only ready} fds. $O(1)$ amortized per ready fd. \\
kqueue & macOS/BSD & Equivalent to epoll. Uses a ``changelist'' API. What Python uses on macOS. \\
io\_uring & Linux 5.1+ & Shared ring buffers between user and kernel space.
           Zero-copy, zero-syscall batching. The future of Linux async I/O. \\
IOCP & Windows & Completion-based (vs readiness-based). Kernel tells you when I/O
       \emph{completed}, not when it's ready. \\
\bottomrule
\end{tabularx}
\end{tcolorbox}
\end{center}

\begin{infobox}
Python's \py{asyncio} uses the \textbf{best available backend} automatically:
\py{epoll} on Linux, \py{kqueue} on macOS/BSD, \py{IOCP} on Windows. You never interact
with these directly --- the event loop abstracts them via the \py{selectors} module.
\end{infobox}

\section{The Event Loop: What Actually Happens}

An event loop is conceptually simple --- an infinite loop:

\begin{lstlisting}[style=python,numbers=none]
while True:
    ready_fds = epoll.wait()        # block until something is ready
    for fd, event in ready_fds:
        callback = registry[fd]
        callback(fd, event)         # resume the coroutine that was waiting
\end{lstlisting}

Every \py{await} in your code is, at the lowest level, registering a file descriptor with the
event loop and yielding control. When the kernel signals that fd is ready, the event loop
resumes the coroutine that was waiting on it.

\begin{warningbox}
This is why blocking the event loop is fatal: if your coroutine runs CPU-heavy code
without yielding, the \texttt{epoll.wait()} call never executes, and \emph{no other
coroutine} can make progress. There is only one thread running the loop.
\end{warningbox}

\section{From NIC to Coroutine: The Full I/O Path}

When data arrives on a network socket, here's the complete path on Linux:

\begin{enumerate}
    \item \textbf{Hard IRQ.} The NIC receives a packet and fires a \emph{hardware interrupt}.
          The CPU stops what it's doing and runs the NIC driver's interrupt handler. This
          handler is minimal --- it acknowledges the interrupt and schedules a soft IRQ.
    \item \textbf{Soft IRQ (NET\_RX\_SOFTIRQ).} The kernel's network stack processes the packet:
          IP reassembly, TCP state machine, checksum validation. Data is placed into the
          \emph{socket receive buffer} (\texttt{sk\_buff}).
    \item \textbf{epoll notification.} The kernel checks if any \texttt{epoll} instance is watching
          this socket's fd. If so, it adds the fd to the epoll's ready list and wakes up the thread
          blocked on \texttt{epoll\_wait()}.
    \item \textbf{Event loop dispatch.} Python's event loop returns from \texttt{epoll\_wait()},
          finds the ready fd, and invokes the callback registered for it.
    \item \textbf{Coroutine resumes.} The callback resumes the coroutine that was suspended at
          an \py{await} expression. Your application code runs until it hits the next \py{await}.
\end{enumerate}

\begin{analogybox}
\begin{itemize}
    \item \textcolor{goblue}{\textbf{Go:}} The same hard IRQ $\to$ soft IRQ $\to$ socket buffer
          chain happens, but Go's \texttt{netpoller} (which uses epoll internally) is integrated
          into the goroutine scheduler. When a goroutine does \texttt{conn.Read()}, the runtime
          registers the fd and parks the goroutine. The fd-ready notification unparks it ---
          same mechanism, different abstraction level.
    \item \textcolor{BrickRed}{\textbf{Rust:}} Tokio's reactor uses \texttt{mio} (a thin wrapper
          over epoll/kqueue). The \texttt{.await} on a \texttt{TcpStream::read()} suspends the
          future. When mio reports readiness, the future is polled again --- identical flow.
\end{itemize}
\end{analogybox}

\section{How Async Database Drivers Work: asyncpg}

You might wonder: ``the database is just sending bytes over TCP --- what makes a driver async
vs sync?'' The answer is entirely in \emph{how it waits for those bytes}.

\subsection{Sync driver: psycopg2}

psycopg2 uses \texttt{libpq} (PostgreSQL's C client library), which does \emph{blocking}
socket I/O internally:

\begin{lstlisting}[style=python,numbers=none]
# Pseudocode of what happens inside psycopg2
sock.send(query_message)     # blocks until sent
response = sock.recv(65536)  # blocks until data arrives (!!!)
return parse(response)
\end{lstlisting}

The \texttt{recv()} call suspends the thread. If you call this from an async function, you've
just blocked the event loop.

\subsection{Async driver: asyncpg}

asyncpg implements the PostgreSQL wire protocol \emph{from scratch} in
Cython --- it doesn't use libpq at all. Instead, it uses non-blocking sockets registered
with the event loop:

\begin{lstlisting}[style=python,numbers=none]
# Pseudocode of what asyncpg does
sock.setblocking(False)
sock.send(query_message)            # non-blocking, returns immediately
# Register fd with event loop and suspend the coroutine
await loop.sock_recv(sock, 65536)   # yields control, resumes on data
return parse(response)
\end{lstlisting}

The key difference: \py{loop.sock\_recv()} registers the socket's file descriptor with epoll,
\emph{yields control back to the event loop}, and only resumes when data arrives. While
waiting for PostgreSQL's response, other coroutines (handling other requests) run freely.

\subsection{PostgreSQL Wire Protocol}

PostgreSQL uses a binary message-based protocol over TCP. A simplified exchange:

\begin{lstlisting}[style=shell,numbers=none]
Client -> Server:  Query('SELECT * FROM users WHERE id = 1')
Server -> Client:  RowDescription(columns=['id','name','email'])
Server -> Client:  DataRow([1, 'Alice', 'alice@example.com'])
Server -> Client:  CommandComplete('SELECT 1')
Server -> Client:  ReadyForQuery(status='idle')
\end{lstlisting}

Each message has a type byte, a 4-byte length, and a payload. asyncpg parses these
incrementally as bytes arrive on the non-blocking socket. The protocol itself is
inherently asynchronous --- the server streams rows as they're produced, and the
client can process them as they arrive.

\begin{infobox}
PostgreSQL also supports an \textbf{extended query protocol} with explicit Parse/Bind/Execute
phases, which asyncpg uses for prepared statements. This enables binary data transfer
(faster than text encoding) and server-side query plan caching. asyncpg's raw performance
advantage over psycopg2 comes largely from this protocol-level optimization, not just from
being non-blocking.
\end{infobox}

\begin{analogybox}
\begin{itemize}
    \item \textcolor{goblue}{\textbf{Go:}} \texttt{pgx} (the standard Go PG driver) uses the same
          wire protocol approach --- it implements the PostgreSQL protocol directly on
          non-blocking sockets, managed by Go's netpoller. Same architecture, different runtime.
    \item \textcolor{BrickRed}{\textbf{Rust:}} \texttt{tokio-postgres} uses the same pattern ---
          async reads/writes on a \texttt{TcpStream} via Tokio's reactor, parsing the PG wire
          protocol in Rust.
    \item \textcolor{javaorange}{\textbf{Java:}} R2DBC PostgreSQL driver (\texttt{r2dbc-postgresql})
          does the same, reimplementing the wire protocol on non-blocking NIO channels instead of
          using JDBC (which wraps libpq-style blocking I/O).
\end{itemize}
\end{analogybox}

% ════════════════════════════════════════════════════════
\chapter{Tooling \& Environment}

Before touching any async code, let's get the toolchain right. Python 3.14 has a modern ecosystem
that may feel unfamiliar if you're coming from older Python or from other languages.

\section{Python 3.14 --- What Changed}

Python 3.14 (released October 2025) is the target runtime. The most impactful changes for backend
developers:

\begin{itemize}
    \item \textbf{Deferred annotation evaluation} (PEP 649) --- annotations are evaluated lazily.
          Forward references work without quoting. If you've used Rust's type system or Kotlin's,
          this will feel natural; Python finally stops choking on circular type references at import time.
    \item \textbf{\py{asyncio.get\_event\_loop()} raises \py{RuntimeError}} --- no more implicit loop
          creation. This is the single biggest breaking change for asyncio code. Think of it like Go
          removing the global default \gocode{context.Background()} and forcing you to pass context explicitly.
    \item \textbf{Template strings (t-strings)} --- tagged template literals similar to JavaScript/Kotlin,
          but designed for safe SQL/HTML interpolation.
    \item \textbf{\py{uuid.uuid7()}} --- time-ordered UUIDs. If you've used ULIDs in other ecosystems,
          this is the stdlib equivalent. Better B-tree index locality than uuid4.
    \item \textbf{Free-threaded build} --- experimental, not for production. Python is slowly moving
          toward removing the GIL, similar to how Go had goroutines from day one.
    \item \textbf{Asyncio introspection} --- \texttt{python -m asyncio ps <PID>} lets you inspect
          running tasks in a live process, analogous to \texttt{goroutine dump} via \texttt{SIGQUIT}
          in Go or \texttt{jstack} in Java.
\end{itemize}

\section{Package Manager: uv}

\begin{dangerbox}
Never use \texttt{pip}, \texttt{poetry}, or \texttt{conda}. The project uses \textbf{uv} exclusively.
\end{dangerbox}

\py{uv} is to Python what \gocode{go mod} is to Go, what \rustcode{cargo} is to Rust, or what
Gradle is to Kotlin --- a unified tool for dependency management and task running. It's written in
Rust and is orders of magnitude faster than pip.

\begin{lstlisting}[style=shell]
# Add a dependency (like `cargo add` or `go get`)
uv add httpx

# Run a tool (like `cargo run` or `go run`)
uv run pytest
uv run ruff check .
\end{lstlisting}

\begin{analogybox}
\begin{tabularx}{\textwidth}{lllll}
\textbf{Action} & \textcolor{goblue}{\textbf{Go}} & \textcolor{BrickRed}{\textbf{Rust}} & \textcolor{javaorange}{\textbf{Java/Kotlin}} & \textcolor{pyblue}{\textbf{Python 3.14}} \\
\midrule
Add dep     & \texttt{go get}          & \texttt{cargo add}   & \texttt{gradle add} & \texttt{uv add} \\
Run         & \texttt{go run}          & \texttt{cargo run}   & \texttt{gradle run} & \texttt{uv run} \\
Lock file   & \texttt{go.sum}          & \texttt{Cargo.lock}  & \texttt{gradle.lock} & \texttt{uv.lock} \\
Manifest    & \texttt{go.mod}          & \texttt{Cargo.toml}  & \texttt{build.gradle} & \texttt{pyproject.toml} \\
\end{tabularx}
\end{analogybox}

\section{Linting, Formatting, Type Checking}

Three tools, each best-in-class:

\begin{itemize}
    \item \textbf{ruff} --- replaces flake8, isort, and black in a single Rust-based binary.
          Think \texttt{rustfmt + clippy} combined, or \texttt{gofmt + go vet} in one tool.
    \item \textbf{ty} (preferred) or \textbf{mypy --strict} --- static type checker.
          \py{ty} is the newer, faster alternative. Analogous to the Go compiler's built-in type
          checking, Rust's \texttt{cargo check}, or Kotlin's compiler.
    \item \textbf{pytest + pytest-asyncio} --- testing framework. pytest is to Python what
          \texttt{go test} is to Go or what JUnit 5 is to Java/Kotlin.
\end{itemize}

% ════════════════════════════════════════════════════════
\chapter{The Event Loop \& Entry Points}

\section{The Single Rule: \py{asyncio.run()}}

If you know one thing, know this: the event loop is not something you ``get'' --- it's something
you \emph{run}. The programming model is:

\begin{enumerate}
    \item Write an \py{async def main()}.
    \item Call \py{asyncio.run(main())} from synchronous code.
    \item Everything async happens inside that call.
\end{enumerate}

\begin{lstlisting}[style=python]
import asyncio

async def main() -> None:
    print("Hello from the event loop")

if __name__ == "__main__":
    asyncio.run(main())
\end{lstlisting}

\begin{analogybox}
\begin{itemize}
    \item \textcolor{goblue}{\textbf{Go:}} You don't create an event loop --- goroutines are
          scheduled by the runtime. \py{asyncio.run()} is roughly equivalent to your \texttt{func main()}
          --- it's the root of all concurrent work.
    \item \textcolor{javaorange}{\textbf{Kotlin:}} This is \javacode{runBlocking \{ ... \}} ---
          it bridges the sync world to the coroutine world.
    \item \textcolor{BrickRed}{\textbf{Rust:}} This is \rustcode{\#[tokio::main]} or
          \texttt{tokio::runtime::Runtime::new().block\_on(main())}.
\end{itemize}
\end{analogybox}

\begin{dangerbox}[title={Removed patterns}]
These will \textbf{crash} in Python 3.14:
\begin{itemize}
    \item \py{asyncio.get\_event\_loop()} --- raises \py{RuntimeError} if no loop is running.
    \item \py{loop.run\_until\_complete()} --- don't use in application code.
    \item \py{asyncio.ensure\_future()} --- removed; use \py{create\_task()}.
    \item \py{@asyncio.coroutine} / \py{yield from} --- deleted generator-based coroutine syntax.
\end{itemize}
\end{dangerbox}

\section{Multi-Phase Startup with \py{asyncio.Runner}}

Sometimes you need to interleave sync and async during startup (e.g., run migrations between
async init phases). Use \py{asyncio.Runner}:

\begin{lstlisting}[style=python]
with asyncio.Runner() as runner:
    runner.run(init_db())       # async phase 1
    apply_migrations()          # sync phase
    runner.run(serve())         # async phase 2
\end{lstlisting}

\begin{infobox}
\py{asyncio.Runner} reuses the same event loop across calls, unlike calling \py{asyncio.run()}
multiple times (which creates and destroys a loop each time). This is useful when you have
resources (connections, caches) that live on the loop.
\end{infobox}

\begin{analogybox}
In Go, this problem doesn't exist because goroutines and sync code coexist naturally. In Rust/tokio,
this is like creating a \rustcode{Runtime} manually and calling \texttt{.block\_on()} multiple times
on the same runtime handle.
\end{analogybox}

% ════════════════════════════════════════════════════════
\chapter{Structured Concurrency}

This is the most important chapter. If you internalize one concept, make it this one.

\section{The Problem with \py{gather()}}

\py{asyncio.gather()} was the old way to run tasks concurrently. It has a critical flaw:
\textbf{when one task fails, the others keep running.} You get a partial failure with zombied
tasks leaking in the background.

\begin{dangerbox}[title={Never use \py{asyncio.gather()}}]
\begin{lstlisting}[style=python,numbers=none]
# WRONG --- siblings continue after failure
results = await asyncio.gather(
    fetch_users(),
    fetch_orders(),
    return_exceptions=True,  # hides errors!
)
\end{lstlisting}
Even with \py{return\_exceptions=True}, you're sweeping failures under the rug. The
remaining tasks run unsupervised.
\end{dangerbox}

\begin{analogybox}
This is the exact problem that motivated:
\begin{itemize}
    \item \textcolor{goblue}{\textbf{Go:}} \gocode{errgroup.Group} from \texttt{golang.org/x/sync/errgroup}
          --- when one goroutine in the group returns an error, the shared context is cancelled,
          signaling all siblings to stop.
    \item \textcolor{javaorange}{\textbf{Java 21+:}} \javacode{StructuredTaskScope} (JEP 453) ---
          \texttt{ShutdownOnFailure} policy cancels siblings on first exception, exactly like \py{TaskGroup}.
    \item \textcolor{BrickRed}{\textbf{Rust:}} \rustcode{tokio::task::JoinSet} --- collects
          tasks and cancels remaining on drop. Or the \texttt{futures::try\_join!} macro that
          cancels on first error.
\end{itemize}
\end{analogybox}

\section{\py{TaskGroup} --- The Right Way}

\py{asyncio.TaskGroup} implements \textbf{structured concurrency}: all tasks spawned within the
group are guaranteed to be finished or cancelled when the \py{async with} block exits.

\begin{lstlisting}[style=python]
async with asyncio.TaskGroup() as tg:
    task_a = tg.create_task(fetch_users(), name="fetch-users")
    task_b = tg.create_task(fetch_orders(), name="fetch-orders")
# Both tasks are done here
users = task_a.result()
orders = task_b.result()
\end{lstlisting}

\begin{successbox}
Key guarantees of \py{TaskGroup}:
\begin{enumerate}
    \item If any task raises, all siblings are \textbf{cancelled}.
    \item The \py{async with} block doesn't exit until \textbf{all tasks} have finished (completed,
          failed, or been cancelled).
    \item Exceptions are collected into an \py{ExceptionGroup}, not silently dropped.
\end{enumerate}
\end{successbox}

\begin{infobox}
The \py{name} parameter on \py{create\_task()} is optional but strongly recommended. It shows up
in tracebacks and in the new \texttt{asyncio ps} debugger, making production debugging
much easier. Think of it like naming goroutines in Go's pprof labels.
\end{infobox}

\section{Exception Groups and \py{except*}}

When multiple tasks fail inside a \py{TaskGroup}, you get an \py{ExceptionGroup} ---
a container that holds multiple exceptions simultaneously. Python 3.11+ introduced
\py{except*} syntax to handle them:

\begin{lstlisting}[style=python]
try:
    async with asyncio.TaskGroup() as tg:
        tg.create_task(might_fail())
        tg.create_task(might_also_fail())
except* ValueError as eg:
    for exc in eg.exceptions:
        log.error("ValueError: %s", exc)
except* ConnectionError as eg:
    for exc in eg.exceptions:
        log.error("Connection failed: %s", exc)
\end{lstlisting}

\begin{warningbox}
\py{except*} is not \py{except}. Key differences:
\begin{itemize}
    \item Multiple \py{except*} clauses can match from the \emph{same} \py{ExceptionGroup}
          --- unlike regular \py{except} where only the first match wins.
    \item \py{eg.exceptions} is a tuple of the matched exceptions within the group.
    \item Unmatched exceptions propagate upward automatically.
\end{itemize}
\end{warningbox}

\begin{analogybox}
\begin{itemize}
    \item \textcolor{goblue}{\textbf{Go:}} Go has no equivalent --- \gocode{errgroup} gives you
          back a single error (the first one). You'd need \texttt{multierr} or a custom
          multi-error type. Python's \py{ExceptionGroup} is richer.
    \item \textcolor{javaorange}{\textbf{Java:}} \javacode{StructuredTaskScope.ShutdownOnFailure}
          also returns only the first exception. The concept of exception groups doesn't exist in
          Java/Kotlin. Python is ahead here.
    \item \textcolor{BrickRed}{\textbf{Rust:}} \rustcode{JoinSet} returns results one-by-one via
          \texttt{.join\_next()}. Multiple errors must be collected manually. No language-level
          multi-error handling.
\end{itemize}
\end{analogybox}

% ════════════════════════════════════════════════════════
\chapter{Timeouts \& Cancellation}

\section{Timeouts with \py{asyncio.timeout()}}

Every external I/O call must be wrapped in a timeout. No exceptions. A hung HTTP call or
database query without a timeout is a resource leak waiting to happen.

\begin{lstlisting}[style=python]
async with asyncio.timeout(5.0):
    data = await client.get("/api/resource")
\end{lstlisting}

For absolute deadlines (e.g., ``finish by timestamp X''):

\begin{lstlisting}[style=python]
deadline = asyncio.get_running_loop().time() + 30.0
async with asyncio.timeout_at(deadline):
    await long_running_operation()
\end{lstlisting}

\begin{dangerbox}[title={Do not use \py{asyncio.wait\_for()}}]
\py{wait\_for()} exists but has subtle cancellation bugs and worse semantics than the
context-manager approach. Always use \py{asyncio.timeout()}.
\end{dangerbox}

\begin{analogybox}
\begin{itemize}
    \item \textcolor{goblue}{\textbf{Go:}} This is \gocode{context.WithTimeout(ctx, 5*time.Second)}.
          The context manager pattern in Python mirrors Go's \texttt{defer cancel()} pattern ---
          when the block exits, the timeout is cleaned up.
    \item \textcolor{javaorange}{\textbf{Kotlin:}} \javacode{withTimeout(5000) \{ ... \}}
          from \texttt{kotlinx.coroutines}. Nearly identical API.
    \item \textcolor{BrickRed}{\textbf{Rust:}} \rustcode{tokio::time::timeout(Duration::from\_secs(5), fut).await}.
          Same idea, function wrapper instead of context manager.
\end{itemize}
\end{analogybox}

\section{Cancellation Discipline}

Cancellation in Python's asyncio works by injecting a \py{CancelledError} at the next
\py{await} point. This is cooperative cancellation --- the coroutine must actually hit an
\py{await} to notice it's been cancelled.

\begin{successbox}[title={The Rules}]
\begin{enumerate}
    \item \textbf{Never swallow \py{CancelledError}}. If you catch it, always re-raise.
    \item Use \py{try/finally} for cleanup, not \py{try/except CancelledError}.
    \item \textbf{Never} call \py{task.cancel()} on tasks managed by a \py{TaskGroup} ---
          the group handles cancellation internally.
\end{enumerate}
\end{successbox}

The cleanup pattern --- \py{try/finally} ensures resources are released even on cancellation:

\begin{lstlisting}[style=python]
async def worker() -> None:
    conn = await connect()
    try:
        await process(conn)
    finally:
        await conn.close()  # runs even on cancellation
\end{lstlisting}

\subsection{How to Initiate Cancellation}

Outside of a \py{TaskGroup}, you cancel a task explicitly with \py{task.cancel()}:

\begin{lstlisting}[style=python]
async def main() -> None:
    task = asyncio.create_task(long_running_work())

    await asyncio.sleep(5.0)
    task.cancel()  # request cancellation

    try:
        await task  # let CancelledError propagate
    except asyncio.CancelledError:
        log.info("task was cancelled cleanly")
\end{lstlisting}

In practice, you rarely call \py{task.cancel()} directly. The higher-level mechanisms do it for you:

\begin{itemize}
    \item \py{asyncio.timeout()} --- cancels on deadline.
    \item \py{TaskGroup} --- cancels siblings when one fails.
\end{itemize}

\begin{analogybox}
\begin{itemize}
    \item \textcolor{goblue}{\textbf{Go:}} Cancellation via \gocode{context.Context} is also
          cooperative --- goroutines check \texttt{ctx.Done()} or \texttt{ctx.Err()}. The difference
          is that Go doesn't inject an exception; you poll. Python injects at \py{await} points.
    \item \textcolor{javaorange}{\textbf{Kotlin:}} Identical model. \javacode{CancellationException}
          is thrown at suspension points. The rule ``never swallow it'' is the same in both languages.
    \item \textcolor{BrickRed}{\textbf{Rust:}} Cancellation happens by \emph{dropping} the future.
          Synchronous \texttt{Drop} destructors \emph{do} run --- RAII cleanup works
          (file handles close, mutex guards release, memory is freed). But any \emph{async}
          cleanup code inside the future body (e.g., \texttt{conn.shutdown().await}) never executes,
          because the future is never polled again.
          For cooperative cancellation with async cleanup, Tokio provides
          \texttt{tokio\_util::sync::CancellationToken} --- the caller sets the token, and the
          task checks it at \texttt{.await} points, similar to Go's context pattern:
\begin{lstlisting}[style=rust,numbers=none]
// Rust: cooperative cancellation with CancellationToken
tokio::select! {
    _ = token.cancelled() => {
        conn.shutdown().await; // async cleanup runs
    }
    result = do_work() => { handle(result); }
}
\end{lstlisting}
\end{itemize}
\end{analogybox}

% ════════════════════════════════════════════════════════
\chapter{Backpressure \& Resource Limits}

\section{Why Backpressure Matters}

Without backpressure, a fast producer can overwhelm a slow consumer, causing unbounded memory
growth and eventual OOM. This is a problem in every async runtime, but Python makes it easy
to forget because queues default to unbounded.

\begin{dangerbox}[title={Unbounded queues are a bug}]
\py{asyncio.Queue()} without \py{maxsize} is unbounded. Always set \py{maxsize}.
\end{dangerbox}

\section{Bounded Queues}

\begin{lstlisting}[style=python]
# Producer blocks (awaits) when queue is full
queue: asyncio.Queue[WorkItem] = asyncio.Queue(maxsize=100)
\end{lstlisting}

\begin{analogybox}
\begin{itemize}
    \item \textcolor{goblue}{\textbf{Go:}} \gocode{make(chan WorkItem, 100)} --- Go channels are
          bounded by default if you specify a capacity. Python queues are unbounded by default.
          Always pass \py{maxsize}.
    \item \textcolor{javaorange}{\textbf{Java:}} \javacode{new ArrayBlockingQueue<>(100)} ---
          Java's bounded blocking queue. Same concept.
    \item \textcolor{BrickRed}{\textbf{Rust:}} \rustcode{tokio::sync::mpsc::channel(100)} ---
          Tokio's bounded channel. The sender's \texttt{.send().await} suspends when full.
\end{itemize}
\end{analogybox}

\section{Semaphores for Concurrency Limiting}

Use \py{asyncio.Semaphore} to limit concurrent access to a shared resource (e.g., connection pool,
rate-limited API):

\begin{lstlisting}[style=python]
sem = asyncio.Semaphore(20)

async def rate_limited_fetch(url: str) -> Response:
    async with sem:
        return await client.get(url)
\end{lstlisting}

\begin{infobox}
For \py{StreamWriter} (raw TCP), always call \py{writer.drain()} after \py{writer.write()}.
The \py{drain()} coroutine suspends until the write buffer is flushed, preventing you from
piling up data faster than the network can send it.
\end{infobox}

% ════════════════════════════════════════════════════════
\chapter{Blocking \& CPU-Bound Work}

\section{The Golden Rule}

\begin{dangerbox}[title={Never block the event loop}]
Running synchronous, blocking, or CPU-intensive code directly in an async function
\textbf{freezes the entire event loop}. No other coroutine can make progress.
This is the \#1 production bug in async Python.
\end{dangerbox}

\begin{analogybox}
\begin{itemize}
    \item \textcolor{goblue}{\textbf{Go:}} Goroutines are preemptively scheduled, so a
          CPU-bound goroutine doesn't block others (since Go 1.14). Python's event loop is
          \emph{cooperative} --- one blocked coroutine stops everything.
    \item \textcolor{javaorange}{\textbf{Kotlin:}} Same problem exists. Kotlin docs warn against
          CPU-bound work in \texttt{Dispatchers.Main}; you use \texttt{Dispatchers.Default} or
          \texttt{Dispatchers.IO} instead. Python's equivalent is \py{to\_thread()} and
          \py{run\_in\_executor()}.
    \item \textcolor{BrickRed}{\textbf{Rust:}} Tokio has the same issue. You use
          \rustcode{tokio::task::spawn\_blocking()} to offload blocking work to a thread pool.
          Python's \py{asyncio.to\_thread()} is the direct equivalent.
\end{itemize}
\end{analogybox}

\section{Offloading Blocking I/O}

For synchronous library calls (legacy DB drivers, file I/O, etc.), use \py{asyncio.to\_thread()}:

\begin{lstlisting}[style=python]
result = await asyncio.to_thread(sync_db_call, query)
\end{lstlisting}

This runs the function in a \py{ThreadPoolExecutor}, keeping the event loop free.

\section{Offloading CPU-Bound Work}

For CPU-intensive work (data processing, image manipulation, cryptography), threads won't help
due to the GIL\footnote{The Global Interpreter Lock prevents multiple threads from executing
Python bytecode simultaneously. The free-threaded build in 3.14 is experimental.}. Use
\py{ProcessPoolExecutor}:

\begin{lstlisting}[style=python]
from concurrent.futures import ProcessPoolExecutor

executor = ProcessPoolExecutor(max_workers=4)

async def heavy_compute(data: bytes) -> Result:
    loop = asyncio.get_running_loop()
    return await loop.run_in_executor(executor, process_data, data)
\end{lstlisting}

\begin{warningbox}
Note the use of \py{asyncio.get\_running\_loop()} here --- not \py{get\_event\_loop()}.
The former returns the loop that's currently running (safe). The latter tries to get or
create a loop (broken in 3.14).
\end{warningbox}

\begin{analogybox}
CPU parallelism varies drastically across languages because of the GIL:
\begin{itemize}
    \item \textcolor{goblue}{\textbf{Go:}} This problem doesn't exist. Goroutines run on multiple
          OS threads by default (\texttt{GOMAXPROCS} = CPU count). CPU-bound goroutines run truly
          in parallel. No GIL, no \texttt{ProcessPoolExecutor} --- just \texttt{go func()}.
    \item \textcolor{javaorange}{\textbf{Java/Kotlin:}} No GIL either. Use
          \javacode{Dispatchers.Default} (sized to CPU cores) for CPU-bound coroutines, or
          \texttt{ForkJoinPool} for parallel streams. Java threads are real OS threads with
          true parallelism.
    \item \textcolor{BrickRed}{\textbf{Rust:}} No GIL. Use \texttt{rayon} for data parallelism
          (parallel iterators) or \rustcode{std::thread::spawn()} for OS threads. In async code,
          \rustcode{tokio::task::spawn\_blocking()} offloads to a thread pool --- but unlike Python,
          the offloaded code runs truly parallel with other threads, not serialized by a lock.
\end{itemize}
\vspace{4pt}
\textbf{Bottom line:} Python is the only language here where CPU-bound work requires process
boundaries for true parallelism. The free-threaded build (3.14, experimental) aims to fix this.
\end{analogybox}

% ════════════════════════════════════════════════════════
\chapter{Type Annotations}

\section{Modern Python Type Syntax}

Python 3.14 uses built-in generics and union syntax. No imports from \py{typing} needed
for basic types.

\begin{migrationbox}
\begin{tabularx}{\textwidth}{lX}
\textbf{Old (pre-3.9)} & \textbf{Python 3.14} \\
\midrule
\py{typing.List[str]} & \py{list[str]} \\
\py{typing.Dict[str, int]} & \py{dict[str, int]} \\
\py{typing.Optional[User]} & \py{User | None} \\
\py{typing.Union[str, bytes]} & \py{str | bytes} \\
\py{typing.Tuple[int, ...]} & \py{tuple[int, ...]} \\
\end{tabularx}
\end{migrationbox}

\begin{lstlisting}[style=python]
# Built-in generics --- no imports needed
def process(items: list[str]) -> dict[str, int]: ...

# Union with | syntax
def parse(value: str | bytes) -> Result: ...

# Async function signatures --- always annotate return type
async def fetch_user(user_id: int) -> User | None: ...

# Abstract types from collections.abc
from collections.abc import (
    AsyncIterator, Callable, Coroutine, Sequence,
)

async def stream_events() -> AsyncIterator[Event]:
    async for event in event_source:
        yield event
\end{lstlisting}

\subsection{Protocol vs ABC: Structural vs Nominal Typing}

The difference is \emph{how} a type satisfies an interface:

\begin{itemize}
    \item \textbf{ABC (nominal):} A class must \emph{explicitly inherit} from the interface.
    \item \textbf{Protocol (structural):} A class satisfies the interface if it has the right
          methods/attributes. No inheritance needed.
\end{itemize}

\begin{lstlisting}[style=python]
# --- ABC approach: requires explicit inheritance ---
from abc import ABC, abstractmethod

class Repository(ABC):
    @abstractmethod
    async def find(self, id: int) -> Model | None: ...

class UserRepo(Repository):  # MUST inherit Repository
    async def find(self, id: int) -> User | None:
        return await db.get(User, id)

# --- Protocol approach: structural matching ---
from typing import Protocol

class Repository(Protocol):
    async def find(self, id: int) -> Model | None: ...

class UserRepo:  # No inheritance! Just has the right shape.
    async def find(self, id: int) -> User | None:
        return await db.get(User, id)

# UserRepo satisfies Repository because it has a matching
# find() method --- the type checker verifies this statically.
async def get_by_id(repo: Repository, id: int) -> Model | None:
    return await repo.find(id)

get_by_id(UserRepo(), 42)  # OK --- structural match
\end{lstlisting}

\begin{successbox}[title={Prefer Protocol}]
Use \py{Protocol} when the interface is consumed by your code (dependency injection,
service boundaries). Use \py{ABC} only when you need to enforce inheritance hierarchies
or share default implementations via methods.
\end{successbox}

\begin{analogybox}
\begin{itemize}
    \item \textcolor{goblue}{\textbf{Go:}} Go interfaces are \emph{always} structural ---
          if a type has the methods, it satisfies the interface. No \texttt{implements} keyword.
          Python's \py{Protocol} is the exact equivalent. Python's \py{ABC} has no Go counterpart.
    \item \textcolor{javaorange}{\textbf{Java/Kotlin:}} Java interfaces are always nominal ---
          you must write \javacode{implements Foo}. This is how Python's \py{ABC} works.
          Java has no structural typing equivalent.
    \item \textcolor{BrickRed}{\textbf{Rust:}} Rust traits are nominal (\rustcode{impl Trait for Type}),
          but trait \emph{bounds} on generics achieve a Protocol-like effect --- you define the
          required shape, and any type that implements those traits is accepted:
\begin{lstlisting}[style=rust,numbers=none]
// Rust: trait bounds constrain by capability
fn get_by_id<R: Repository>(repo: &R, id: i64) -> Option<Model> {
    repo.find(id)    // R must impl Repository
}
// Any type that `impl Repository` works --- explicit impl,
// but generic bounds give the same "shape matching" feel.

// Trait objects for dynamic dispatch (like Python's Protocol):
async fn get_by_id(repo: &dyn Repository, id: i64) -> Option<Model>
\end{lstlisting}
\end{itemize}
\end{analogybox}

\section{Deferred Annotation Evaluation (PEP 649)}

In Python 3.14, annotations are evaluated lazily. This means forward references work without
quoting:

\begin{lstlisting}[style=python]
# This works in 3.14 --- no need for quotes or
# from __future__ import annotations
class TreeNode:
    def add_child(self, child: TreeNode) -> None: ...
\end{lstlisting}

\begin{infobox}
In older Python, you'd need \py{"TreeNode"} (string annotation) or
\py{from \_\_future\_\_ import annotations}. Python 3.14 makes this unnecessary by
evaluating annotations only when explicitly requested (e.g., by type checkers or
\py{get\_type\_hints()}).
\end{infobox}

% ════════════════════════════════════════════════════════
\chapter{Async Library Choices}

Choosing the right async library matters. The wrong choice means fighting the framework
instead of building your application.

\section{Recommended Stack}

\begin{center}
\begin{tcolorbox}[
    colback=codebg,
    colframe=pyblue,
    width=0.95\textwidth,
    arc=3pt,
    boxrule=1pt,
]
\renewcommand{\arraystretch}{1.4}
\begin{tabularx}{\textwidth}{>{\bfseries}l l X}
\toprule
Purpose & Library & Notes \\
\midrule
HTTP client & \py{httpx} & Async-native, HTTP/2, sync+async in one library \\
Web framework & \py{FastAPI} & Async-first, Pydantic v2 validation \\
ORM / DB & \py{SQLAlchemy 2.x} async & \texttt{from sqlalchemy.ext.asyncio import ...} \\
PostgreSQL & \py{asyncpg} & Fastest async PG driver \\
Redis & \py{redis.asyncio} & Official async support \\
WebSocket & \py{websockets} & Or FastAPI's built-in WebSocket support \\
Task queue & \py{arq} or \py{taskiq} & Async-native alternatives to Celery \\
\bottomrule
\end{tabularx}
\end{tcolorbox}
\end{center}

\begin{dangerbox}[title={Do not use aiohttp}]
For new projects, use \py{httpx.AsyncClient} instead of \py{aiohttp}. httpx has a better API,
supports HTTP/2, and provides both sync and async clients in one library (useful for testing).
\end{dangerbox}

\begin{analogybox}
\begin{itemize}
    \item \textcolor{goblue}{\textbf{Go:}} Go's \texttt{net/http} is both sync and async
          transparently. httpx is the closest Python equivalent --- one library, both modes.
    \item \textcolor{javaorange}{\textbf{Java/Kotlin:}} FastAPI is comparable to Ktor (Kotlin)
          or Spring WebFlux (Java) --- async-first web framework with validation.
    \item \textcolor{BrickRed}{\textbf{Rust:}} httpx $\approx$ \texttt{reqwest}, FastAPI $\approx$
          \texttt{axum}, SQLAlchemy async $\approx$ \texttt{sqlx}.
\end{itemize}
\end{analogybox}

% ════════════════════════════════════════════════════════
\chapter{Project Structure}

\section{Canonical Layout}

\begin{lstlisting}[style=directory]
project/
|-- pyproject.toml
|-- CLAUDE.md
|-- src/
|   +-- myapp/
|       |-- __init__.py
|       |-- main.py          # asyncio.run() entry
|       |-- config.py        # pydantic Settings
|       |-- models/          # domain models / SQLAlchemy
|       |-- services/        # business logic (async)
|       |-- api/             # FastAPI routers
|       |-- clients/         # external HTTP/gRPC clients
|       +-- infra/           # DB, cache, queue adapters
|-- tests/
|   |-- conftest.py          # async fixtures
|   |-- unit/
|   +-- integration/
+-- uv.lock
\end{lstlisting}

\begin{analogybox}
\begin{itemize}
    \item \textcolor{goblue}{\textbf{Go:}} The \texttt{src/myapp/} structure with sub-packages
          mirrors Go's \texttt{internal/} package pattern. \texttt{main.py} is your \texttt{cmd/server/main.go}.
    \item \textcolor{BrickRed}{\textbf{Rust:}} \texttt{src/myapp/} $\approx$ \texttt{src/},
          \texttt{main.py} $\approx$ \texttt{main.rs}. The \texttt{services/} and \texttt{infra/}
          split mirrors the hexagonal architecture common in Rust web apps.
    \item \textcolor{javaorange}{\textbf{Java/Kotlin:}} This follows the same layered architecture:
          \texttt{api/} = controllers, \texttt{services/} = service layer, \texttt{infra/} = repositories.
\end{itemize}
\end{analogybox}

\begin{infobox}
Key files:
\begin{itemize}
    \item \textbf{\texttt{pyproject.toml}} --- the single source of truth for project metadata,
          dependencies, and tool configuration (pytest, ruff, ty).
    \item \textbf{\texttt{uv.lock}} --- deterministic lockfile. Commit it. Never edit manually.
    \item \textbf{\texttt{conftest.py}} --- shared pytest fixtures. Async fixtures go here.
\end{itemize}
\end{infobox}

% ════════════════════════════════════════════════════════
\chapter{Testing Async Code}

\section{Setup: pytest-asyncio with Auto Mode}

Add this to \texttt{pyproject.toml}:

\begin{lstlisting}[style=toml]
[tool.pytest.ini_options]
asyncio_mode = "auto"
\end{lstlisting}

With \py{auto} mode, any \py{async def test\_...} function is automatically recognized as
an async test. No decorators needed.

\begin{migrationbox}
Without auto mode, you'd need \py{@pytest.mark.asyncio} on every async test. Auto mode
eliminates this boilerplate entirely.
\end{migrationbox}

\section{Writing Async Tests}

\begin{lstlisting}[style=python]
# No decorator needed with mode=auto
async def test_fetch_user(db_session: AsyncSession) -> None:
    user = await create_user(db_session, name="Alex")
    fetched = await get_user(db_session, user.id)
    assert fetched is not None
    assert fetched.name == "Alex"
\end{lstlisting}

\section{Async Fixtures}

Fixtures can be async too. Use \py{AsyncIterator} for fixtures that need cleanup:

\begin{lstlisting}[style=python]
from collections.abc import AsyncIterator

@pytest.fixture
async def db_session() -> AsyncIterator[AsyncSession]:
    async with async_session_maker() as session:
        yield session
        await session.rollback()
\end{lstlisting}

\begin{analogybox}
\begin{itemize}
    \item \textcolor{goblue}{\textbf{Go:}} Go's \texttt{testing.T} with \texttt{t.Cleanup()}
          is the rough equivalent. Python's fixture system is more powerful --- fixtures are
          dependency-injected by name.
    \item \textcolor{javaorange}{\textbf{Java/Kotlin:}} Think JUnit 5's \texttt{@BeforeEach}
          and \texttt{@AfterEach}, but with dependency injection and composability.
    \item \textcolor{BrickRed}{\textbf{Rust:}} Rust doesn't have built-in fixtures. Libraries
          like \texttt{rstest} provide similar DI-based fixture injection.
\end{itemize}
\end{analogybox}

\section{Testing TaskGroup Error Handling}

Use \py{pytest.raises(ExceptionGroup)} to test that \py{TaskGroup} failures propagate correctly:

\begin{lstlisting}[style=python]
async def test_task_failure_cancels_siblings() -> None:
    results: list[str] = []
    with pytest.raises(ExceptionGroup) as exc_info:
        async with asyncio.TaskGroup() as tg:
            tg.create_task(append_after(results, "ok", 0.1))
            tg.create_task(fail_after(0.05))
    assert len(exc_info.value.exceptions) == 1
\end{lstlisting}

% ════════════════════════════════════════════════════════
\chapter{Error Handling}

\section{Application Error Hierarchy}

Define a base exception class and derive specific errors from it. Keep it simple:

\begin{lstlisting}[style=python]
class AppError(Exception):
    """Base for all application errors."""

class NotFoundError(AppError):
    """Resource not found."""

class ValidationError(AppError):
    """Input validation failed."""
\end{lstlisting}

\begin{successbox}[title={Error propagation strategy}]
\begin{enumerate}
    \item \textbf{Service layer:} raise errors, don't catch them.
    \item \textbf{API boundary:} catch and convert to HTTP responses.
    \item \textbf{Let it crash:} unexpected errors should propagate to the top and be logged.
\end{enumerate}
\end{successbox}

\begin{lstlisting}[style=python]
# Service layer --- raise, don't catch
async def get_user(user_id: int) -> User:
    user = await repo.find(user_id)
    if user is None:
        raise NotFoundError(f"User {user_id} not found")
    return user

# API boundary --- convert to HTTP
@app.exception_handler(NotFoundError)
async def not_found_handler(
    request: Request, exc: NotFoundError,
) -> JSONResponse:
    return JSONResponse(
        status_code=404,
        content={"detail": str(exc)},
    )
\end{lstlisting}

\begin{analogybox}
\begin{itemize}
    \item \textcolor{goblue}{\textbf{Go:}} Go uses error values, not exceptions. The pattern of
          ``propagate errors up, handle at boundaries'' is the same, but expressed with
          \texttt{if err != nil \{ return err \}} instead of try/except.
    \item \textcolor{javaorange}{\textbf{Java/Kotlin:}} This mirrors Spring's
          \texttt{@ExceptionHandler} pattern. Service throws, controller translates to HTTP status.
    \item \textcolor{BrickRed}{\textbf{Rust:}} This is the \rustcode{Result<T, E>} pattern with
          the \texttt{?} operator for propagation and an error handler middleware in axum/actix.
\end{itemize}
\end{analogybox}

% ════════════════════════════════════════════════════════
\chapter{Graceful Shutdown}

\section{Signal Handling in Async Services}

Production services must handle SIGTERM (from container orchestrators) and SIGINT (Ctrl+C)
cleanly. The pattern:

\begin{lstlisting}[style=python]
import asyncio
import signal

async def main() -> None:
    shutdown_event = asyncio.Event()
    loop = asyncio.get_running_loop()

    for sig in (signal.SIGTERM, signal.SIGINT):
        loop.add_signal_handler(sig, shutdown_event.set)

    server = await start_server()
    await shutdown_event.wait()

    # Graceful shutdown
    server.close()
    await server.wait_closed()
    await cleanup_resources()
\end{lstlisting}

\begin{infobox}
\py{asyncio.Event} is used as a signal bridge between the OS signal (sync) and the async world.
When the signal fires, it sets the event, which unblocks \py{shutdown\_event.wait()}.
\end{infobox}

\begin{analogybox}
\begin{itemize}
    \item \textcolor{goblue}{\textbf{Go:}} This is the
          \gocode{signal.NotifyContext(ctx, ...)} pattern. Context cancellation triggers cleanup
          via \texttt{defer} and \texttt{<-ctx.Done()}.
    \item \textcolor{javaorange}{\textbf{Java:}} \javacode{Runtime.getRuntime().addShutdownHook(...)}.
          Python's approach is more integrated with the async runtime.
    \item \textcolor{BrickRed}{\textbf{Rust:}} \rustcode{tokio::signal::ctrl\_c().await}
          or \texttt{signal::unix::signal(...)}. Same cooperative pattern.
\end{itemize}
\end{analogybox}

% ════════════════════════════════════════════════════════
\chapter{Logging}

\section{Structured Logging Only}

Use \py{structlog} or stdlib \py{logging} with a JSON formatter. \textbf{Never use \py{print()}
for application output.}

\begin{lstlisting}[style=python]
import logging

logger = logging.getLogger(__name__)

# In async code --- log with context
logger.info(
    "processing request",
    extra={"user_id": user_id, "endpoint": path},
)
\end{lstlisting}

\begin{warningbox}
Why not \py{print()}? It's not structured (can't be parsed by log aggregators), has no log levels,
no timestamps, and goes to stdout without buffering control. In production, unstructured logs
are nearly useless for debugging.
\end{warningbox}

\begin{analogybox}
\begin{itemize}
    \item \textcolor{goblue}{\textbf{Go:}} This is equivalent to using \texttt{slog} (structured
          logging, added in Go 1.21) instead of \texttt{fmt.Println}. Same motivation.
    \item \textcolor{javaorange}{\textbf{Java/Kotlin:}} SLF4J + Logback with JSON encoder.
          \py{structlog} serves the same role.
    \item \textcolor{BrickRed}{\textbf{Rust:}} \texttt{tracing} crate with \texttt{tracing-subscriber}.
          Same structured, context-carrying approach.
\end{itemize}
\end{analogybox}

% ════════════════════════════════════════════════════════
\chapter{Python 3.14 Specific Features}

A reference of Python 3.14 features you should actively use in this project.

\section{Quick Reference}

\begin{center}
\renewcommand{\arraystretch}{1.4}
\begin{tcolorbox}[
    colback=codebg,
    colframe=pyblue,
    width=0.98\textwidth,
    arc=3pt,
    boxrule=1pt,
]
\begin{tabularx}{\textwidth}{>{\bfseries\ttfamily}l X}
\toprule
Feature & When to Use \\
\midrule
uuid.uuid7() & New DB primary keys --- time-ordered, better index perf than uuid4 \\
asyncio ps/pstree & Production debugging of stuck/hanging tasks \\
t-strings & Safe SQL/HTML templating (library support required) \\
date.strptime() & Parse dates directly instead of going through datetime \\
Protocol & Interfaces --- prefer over ABC for structural typing \\
PEP 649 & Forward references work without quotes --- just write types \\
\bottomrule
\end{tabularx}
\end{tcolorbox}
\end{center}

\section{\py{uuid.uuid7()} --- Time-Ordered UUIDs}

\begin{lstlisting}[style=python]
import uuid

# Old way --- random, poor index locality
pk = uuid.uuid4()

# New way --- time-ordered, better B-tree performance
pk = uuid.uuid7()
\end{lstlisting}

\begin{infobox}
UUIDv7 encodes a millisecond timestamp in the most significant bits. This means that
UUIDs generated close in time are close in sort order, which dramatically improves B-tree
index performance for insert-heavy workloads. If you've used ULIDs, KSUIDs, or Snowflake IDs
in other ecosystems, this is the stdlib equivalent.
\end{infobox}

\section{Asyncio Process Introspection}

Debug stuck async tasks in a running production process:

\begin{lstlisting}[style=shell]
# List all running tasks (like `goroutine dump` in Go)
python -m asyncio ps <PID>

# Show task parent-child tree
python -m asyncio pstree <PID>
\end{lstlisting}

\begin{analogybox}
\begin{itemize}
    \item \textcolor{goblue}{\textbf{Go:}} \texttt{SIGQUIT} or \texttt{runtime/pprof} goroutine dump
    \item \textcolor{javaorange}{\textbf{Java:}} \texttt{jstack <PID>} or \texttt{Thread.getAllStackTraces()}
    \item \textcolor{BrickRed}{\textbf{Rust:}} \texttt{tokio-console} for inspecting Tokio tasks
\end{itemize}
\end{analogybox}

% ════════════════════════════════════════════════════════
\chapter{The Anti-Patterns List}

A comprehensive list of things that will cause bugs, crashes, or maintenance headaches.
Each item is a hard rule, not a suggestion.

\vspace{12pt}

\begin{center}
\renewcommand{\arraystretch}{1.5}
\begin{tcolorbox}[
    colback=dangerbg,
    colframe=dangerborder,
    coltitle=white,
    colbacktitle=dangerborder,
    fonttitle=\bfseries\sffamily\large,
    title={Banned Patterns},
    arc=3pt,
    boxrule=1.5pt,
    width=0.98\textwidth,
]
\small
\begin{tabularx}{\textwidth}{>{\ttfamily\color{dangerborder}}p{4.5cm} X}
\toprule
\normalfont\textbf{Do NOT use} & \textbf{Use instead} \\
\midrule
asyncio.gather()           & \py{asyncio.TaskGroup} \\
asyncio.get\_event\_loop() & \py{asyncio.run()} or \py{asyncio.get\_running\_loop()} \\
asyncio.ensure\_future()   & \py{tg.create\_task()} within a TaskGroup \\
@asyncio.coroutine         & \py{async def} (generator syntax was removed) \\
yield from                 & \py{await} \\
loop.run\_until\_complete() & \py{asyncio.run()} \\
asyncio.wait\_for()        & \py{asyncio.timeout()} context manager \\
aiohttp                    & \py{httpx.AsyncClient} \\
pip install                & \py{uv add} \\
typing.List / Dict / Optional & \py{list} / \py{dict} / \py{X | None} \\
bare create\_task()         & \py{tg.create\_task()} inside a TaskGroup \\
blocking code in async     & \py{asyncio.to\_thread()} or \py{run\_in\_executor()} \\
catch CancelledError silently & always re-raise after cleanup \\
print()                    & \py{logging} or \py{structlog} \\
\bottomrule
\end{tabularx}
\end{tcolorbox}
\end{center}

% ════════════════════════════════════════════════════════
\chapter*{Appendix: Mental Model Map}
\addcontentsline{toc}{chapter}{Appendix: Mental Model Map}

If you're coming from another language, use this table to map concepts:

\vspace{12pt}

\begin{center}
\renewcommand{\arraystretch}{1.4}
\begin{tcolorbox}[
    colback=white,
    colframe=pyblue!50!black,
    coltitle=white,
    colbacktitle=pyblue!70!black,
    fonttitle=\bfseries\sffamily\large,
    title={Concept Translation Table},
    arc=3pt,
    boxrule=1.5pt,
    width=0.98\textwidth,
]
\scriptsize
\setlength{\tabcolsep}{3pt}
\newcommand{\tcell}[1]{\raggedright\arraybackslash #1}
\begin{tabularx}{\textwidth}{>{\bfseries\raggedright\arraybackslash}p{1.8cm}
    >{\raggedright\arraybackslash}p{2.5cm}
    >{\raggedright\arraybackslash}p{2.2cm}
    >{\raggedright\arraybackslash}p{2.5cm}
    >{\raggedright\arraybackslash}p{2.5cm}}
\toprule
Concept &
\textcolor{pyblue}{\textbf{Python 3.14}} &
\textcolor{goblue}{\textbf{Go}} &
\textcolor{javaorange}{\textbf{Kotlin}} &
\textcolor{BrickRed}{\textbf{Rust/Tokio}} \\
\midrule
Entry point &
\py{asyncio.run()} &
\gocode{func main()} &
\javacode{runBlocking} &
\rustcode{\#[tokio::main]} \\
Spawn task &
\py{tg.create\_task()} &
\gocode{go func()} &
\javacode{launch \{\}} &
\rustcode{tokio::spawn()} \\
Structured group &
\py{TaskGroup} &
\gocode{errgroup} &
\javacode{TaskScope} &
\rustcode{JoinSet} \\
Timeout &
\py{asyncio.timeout()} &
\gocode{ctx.WithTimeout} &
\javacode{withTimeout} &
\rustcode{time::timeout} \\
Cancellation &
\py{CancelledError} &
\gocode{ctx.Err()} &
\javacode{Cancellation\-Exception} &
Drop future \\
Bounded queue &
\py{Queue(maxsize=N)} &
\gocode{make(chan, N)} &
\javacode{Channel(N)} &
\rustcode{mpsc::channel(N)} \\
Offload blocking &
\py{to\_thread()} &
N/A (preemptive) &
\javacode{Dispatchers.IO} &
\rustcode{spawn\_blocking} \\
Semaphore &
\py{asyncio.Semaphore} &
\gocode{semaphore.\-Weighted} &
\javacode{Semaphore} &
\rustcode{sync::Semaphore} \\
HTTP client &
\py{httpx} &
\gocode{net/http} &
\javacode{Ktor / OkHttp} &
\rustcode{reqwest} \\
Web framework &
\py{FastAPI} &
\gocode{net/http} &
\javacode{Ktor} &
\rustcode{axum} \\
Package mgr &
\py{uv} &
\gocode{go mod} &
\javacode{Gradle} &
\rustcode{cargo} \\
Type checking &
\py{ty / mypy} &
Compiler &
Compiler &
Compiler \\
Interface &
\py{Protocol} &
\gocode{interface} &
\javacode{interface} &
\rustcode{trait} \\
\bottomrule
\end{tabularx}
\end{tcolorbox}
\end{center}

\vfill

\begin{center}
\rule{0.4\textwidth}{0.5pt}\\[6pt]
{\small\sffamily\color{Gray}
Generated from project CLAUDE.md --- March 2026}
\end{center}

\end{document}
